\begin{table}[!t]
\centering
\caption{Standardized Effect Sizes}
\label{tab:sde}
\begin{adjustbox}{max width=\textwidth}
\begin{tabular}{lcccccc}
\hline\hline
Outcome & $\hat{\beta}$ & SE & SD($Y$) & SDE & SE(SDE) & Classification \\
\hline
\multicolumn{7}{l}{\textit{Panel A: Pooled}} \\[3pt]
Log total enterprises & -0.008 & 0.045 & 0.114 & -0.066 & 0.393 & Moderate negative \\
Large enterprise share & -0.0004 & 0.0006 & 0.0003 & -1.073 & 1.845 & Large negative \\
Threshold ratio & 0.016 & 0.145 & 0.022 & 0.731 & 6.520 & Large positive \\
Log large (250+) enterprises & -0.340 & 0.256 & 0.205 & -1.656 & 1.248 & Large negative \\[6pt]
\multicolumn{7}{l}{\textit{Panel B: Heterogeneous (sample splits)}} \\[3pt]
England only (food vs.\ controls) & 0.115 & 0.077 & 0.114 & 1.011 & 0.675 & Large positive \\
Scotland/Wales (food vs.\ controls) & 0.123 & 0.049 & 0.243 & 0.505 & 0.203 & Large positive \\
\hline\hline
\end{tabular}
\end{adjustbox}
\begin{tablenotes}[flushleft]
\small
\item \textit{Notes:} \textbf{Country:} United Kingdom (England treated; Scotland and Wales as controls). \textbf{Research question:} Does mandatory calorie labeling for food businesses with 250 or more employees alter the size distribution or total count of food service enterprises? \textbf{Policy mechanism:} The Calorie Labelling (Out of Home Sector) (England) Regulations 2022 require food businesses with 250+ employees to display calorie information on menus, online platforms, and food delivery apps; businesses below 250 employees are exempt, creating a compliance threshold. \textbf{Outcome definition:} Log total enterprise count, share of large enterprises, and threshold ratio (250--499 / 100--249 band) from ONS UK Business Counts. \textbf{Treatment:} Binary (England post-April 2022 vs.\ Scotland/Wales and pre-period). \textbf{Data:} ONS UK Business Counts (NOMIS NM\_142\_1), annual March snapshots, 2010--2024, country $\times$ industry level, $N = 225$. \textbf{Method:} Triple-difference (country $\times$ industry $\times$ time) with two-way clustered standard errors at the unit level; permutation inference ($p = 0.935$). \textbf{Sample:} Food and beverage service enterprises (SIC 56) in England, Scotland, and Wales; control sectors are retail (47), accommodation (55), IT services (62), and legal/accounting (69). SDE $= \hat{\beta} / \text{SD}(Y)$ where SD($Y$) is the pre-treatment standard deviation. Classification refers to magnitude, not statistical significance: Large ($|$SDE$| > 0.15$), Moderate ($0.05$--$0.15$), Small ($0.005$--$0.05$), Null ($< 0.005$).
\end{tablenotes}
\end{table}
